\section{强化学习训练框架}

\subsection{框架演进：从TRL到现代框架}

崔干曲作为RL训练框架的"资深用户"，见证了从TRL到OpenRLHF再到veRL、SLIME的演进历程。他的核心观察是：

\begin{importantbox}{训练框架的关键进化}
\textbf{OpenRLHF的历史性贡献}：率先引入了"高效推理引擎（vLLM）+ 训练引擎（DeepSpeed）解耦并拼接"的设计范式。体验过TRL初期训练效率的人会深刻感受到，这是一个代际跃迁。此后，几乎所有主流框架都采用了这种推理-训练解耦的架构。
\end{importantbox}

\subsection{OpenRLHF的设计哲学}

胡建分享了OpenRLHF的设计初衷：当时作为用户，没有一个框架能让RL工程师用得"爽和顺手"。设计的核心思路是：

\begin{enumerate}
    \item \textbf{从RL工程师的视角出发}：不是System Engineer的视角，而是"以前做PPO跑游戏环境的人"的视角
    \item \textbf{技术选型}：Ray（分布式调度）+ vLLM（推理引擎）+ DeepSpeed（训练引擎）的组合
    \item \textbf{易用性优先}：在保证性能够用的前提下，持续优化易用性
\end{enumerate}

\begin{knowledgebox}{DeepSpeed到Megatron-Core的演进}
胡建解释了为何OpenRLHF最初选择DeepSpeed而非Megatron-Core：当时的Megatron-Core对RL工程师极不友好。后来NVIDIA工程师借鉴了DeepSpeed的优势进行改进，使得Megatron-Core在后来的veRL等框架中变得可用。技术选型需要因时制宜——选择当前最适合目标用户群体的方案。
\end{knowledgebox}

\subsection{理想训练框架的特征}

综合四位嘉宾的观点，理想的RL训练框架应具备以下特征：

\begin{importantbox}{好的RL训练框架}
\begin{enumerate}
    \item \textbf{简洁性}：框架不应过于臃肿，Data Flow清晰，方便Debug
    \item \textbf{可Debug性}：Rollout、Training、Recompute、Reward/Advantage计算等每一步都能打出中间结果并存储——这是正确性的保障
    \item \textbf{易用性}：新手能快速上手，算法研究者能方便地修改调度流程和算法逻辑
    \item \textbf{借助开源力量}：使用开源的推理引擎（vLLM/SGLang）和训练引擎（Megatron），而非自己从零搭建——开源社区的测试覆盖和迭代速度远超内部团队
    \item \textbf{社区生态}：活跃的Contributor、良好的讨论氛围、持续的维护——这决定了框架的长期生命力和影响力
\end{enumerate}
\end{importantbox}

\begin{warningbox}{当前框架的痛点}
\begin{itemize}
    \item veRL等框架逐渐变得臃肿，当需要对调度流程进行修改时"非常难受"
    \item 正确性问题——有团队使用正确性不满足的框架训练，导致模型训不上去且找不到原因
    \item 大规模训练（600B+模型）需要商业公司内部Infra团队做Megatron-Core的定制优化，开源框架难以覆盖
\end{itemize}
\end{warningbox}

\subsection{框架的未来方向}

李英儒从强化学习特性出发，指出了框架需要支持的新需求：

\begin{enumerate}
    \item \textbf{异步训练}：Rollout与Training的分离式架构，支持异步数据流
    \item \textbf{Replay Buffer}：支持数据复用以提升Sample Efficiency
    \item \textbf{可扩展的Rollout}：通过增加Rollout节点来线性扩展采样能力
    \item \textbf{长程Agentic支持}：Agent与Sandbox的长时间交互（coding环境数十小时交互），对框架的调度和资源管理提出新挑战
\end{enumerate}

\begin{knowledgebox}{从游戏RL到大模型RL的框架演进}
李英儒指出，Agentic RL的框架需求越来越像传统游戏RL——Agent跟环境交互、异步采样、分离式架构。但新的挑战在于：除了GPU上的Generation Bottleneck，还增加了\textbf{环境交互Bottleneck}（如等待代码编译、Sandbox反馈等），这使得Sample Efficiency问题更加突出。
\end{knowledgebox}

\subsection{本章小结}

RL训练框架经历了从TRL到OpenRLHF再到veRL/SLIME的快速演进。推理-训练解耦是里程碑式的架构创新。好的框架需要在简洁性、可Debug性、易用性之间取得平衡，同时借助开源社区的力量保持迭代速度。未来的框架需要支持异步训练、长程Agentic交互等新场景。


